\subsection{Divergence, Upper and Lower Limits, and Sequence Estimates}

\begin{definition} \label{def:sequence-diverges-to-infinity} \hlt{Divergence to $\pm\infty$}\\
Let $(s_n)$ be a real sequence.
\begin{enumerate}[label=(\roman*)]
\setlength{\itemsep}{0pt}
\item $s_n\to+\infty\Longleftrightarrow(\forall M\in\R)(\exists N\in\N)(\forall n\geq N)\,s_n\geq M$.
\item $s_n\to-\infty\Longleftrightarrow(\forall M\in\R)(\exists N\in\N)(\forall n\geq N)\,s_n\leq M$.
\end{enumerate}
\end{definition}

\begin{proposition} \label{prop:monotone-sequence-extended-limit} \hlt{Monotone Sequences Have Extended Limits}\\
Let $(s_n)$ be a real sequence.
\begin{enumerate}[label=(\roman*)]
\setlength{\itemsep}{0pt}
\item If $(s_n)$ is increasing, then $s_n\to\sup\{s_n:n\in\N\}$ in $\overline{\R}$.
\item If $(s_n)$ is decreasing, then $s_n\to\inf\{s_n:n\in\N\}$ in $\overline{\R}$.
\item In either case, the extended limit lies in $\R$ iff $(s_n)$ is bounded.
\end{enumerate}
\begin{proof}
\begin{enumerate}[label=(\roman*)]
\setlength{\itemsep}{0pt}
\item[$(i)$] Set $a\defeq\sup\{s_n:n\in\N\}\in\overline{\R}$. Since the range is nonempty, $a\neq-\infty$.\\
If $a\in\R$, fix $\varepsilon>0$ and choose $N\in\N$ with $a-\varepsilon<s_N\leq a$ by the defining property of $a$.\\
Then $(\forall n\geq N)\,a-\varepsilon<s_n\leq a$ by monotonicity, so $s_n\to a$.\\
If $a=+\infty$, fix $M\in\R$ and choose $N\in\N$ with $s_N\geq M$ by the defining property of $a$.\\
Then $(\forall n\geq N)\,s_n\geq M$ by monotonicity, so $s_n\to+\infty$.
\item[$(ii)$] Set $b\defeq\inf\{s_n:n\in\N\}\in\overline{\R}$. Since the range is nonempty, $b\neq+\infty$.\\
If $b\in\R$, fix $\varepsilon>0$ and choose $N\in\N$ with $b\leq s_N<b+\varepsilon$ by the defining property of $b$.\\
Then $(\forall n\geq N)\,b\leq s_n<b+\varepsilon$ by monotonicity, so $s_n\to b$.\\
If $b=-\infty$, fix $M\in\R$ and choose $N\in\N$ with $s_N\leq M$ by the defining property of $b$.\\
Then $(\forall n\geq N)\,s_n\leq M$ by monotonicity, so $s_n\to-\infty$.
\item[$(iii)$] ($\Rightarrow$) In the increasing case, a finite limit $a$ bounds the range above, while $s_1$ bounds it below.\\
In the decreasing case, a finite limit $b$ bounds the range below, while $s_1$ bounds it above. Thus the range is bounded.
\item[$(iii)$] ($\Leftarrow$) If the range is bounded, its supremum in (i) or infimum in (ii) is real by the extended-extremum definitions.
\end{enumerate}
\hlt{Structural Commentary}: Extended extrema split the finite $\varepsilon$-trap and infinite threshold cases.\\
\hlt{Key Insight}: Monotonicity propagates one extremal witness through the entire tail.
\end{proof}
\end{proposition}

\begin{proposition} \label{prop:unbounded-real-sequence-has-infinite-subsequence} \hlt{Unbounded Real Sequences Have Infinite Subsequences}\\
A real sequence is unbounded iff some subsequence tends to $+\infty$ or $-\infty$.
\begin{proof}
\begin{enumerate}[label=\roman*.]
\setlength{\itemsep}{0pt}
\item[$(\Rightarrow)$] An unbounded real sequence is unbounded above or below. If it is unbounded above, every tail is unbounded above.\\
Set $n_0\defeq0$ and recursively choose $n_k>n_{k-1}$ with $s_{n_k}\geq k$. For $M\in\R$, choose $K\in\N$ with $K\geq M$ by Theorem \ref{thm:archimedean-and-rational-density}(i).\\
Then $(\forall k\geq K)\,s_{n_k}\geq k\geq M$, so $s_{n_k}\to+\infty$. If $(s_n)$ is unbounded below, same construction with $s_{n_k}\leq-k$ satisfies $s_{n_k}\to-\infty$.
\item[$(\Leftarrow)$] If $s_{n_k}\to+\infty$, then $(\forall M\in\R)(\exists k\in\N)\,s_{n_k}\geq M$, so $(s_n)$ is unbounded above.\\
If $s_{n_k}\to-\infty$, then $(\forall M\in\R)(\exists k\in\N)\,s_{n_k}\leq M$, so $(s_n)$ is unbounded below.
\end{enumerate}
\hlt{Structural Commentary}: Tail unboundedness supports a recursive threshold subsequence with levels $\pm k$.\\
\hlt{Key Insight}: Increasing indices convert unbounded tails into extended divergence.
\end{proof}
\end{proposition}

\begin{illustrationfigure}[H]
  \centering
  \begin{tikzpicture}[
  >=Stealth,
  x=0.9cm,
  y=0.9cm,
  every node/.style={font=\small},
  axis/.style={->, gray!70, line width=0.55pt},
  zerodot/.style={circle, fill=orange!85!black, inner sep=1.45pt},
  highdot/.style={circle, fill=blue!70!black, inner sep=1.45pt},
  guide/.style={dashed, line width=0.55pt}
]
  \draw[axis] (0.55,0) -- (7.7,0) node[right] {$n$};
  \draw[axis] (0.80,-0.16) -- (0.80,3.55) node[above] {$s_n$};

  \foreach \x in {1.20,2.55,3.90,5.25,6.60} {
    \node[zerodot] at (\x,0) {};
  }
  \foreach \x/\y in {1.85/0.80,3.20/1.35,4.55/1.90,5.90/2.45,7.25/3.00} {
    \node[highdot] at (\x,\y) {};
  }

  \draw[guide, blue!60!black] (1.85,0.80) -- (7.25,3.00);
  \draw[guide, orange!85!black] (0.95,0) -- (6.85,0);
\end{tikzpicture}

  \caption{For $s_n=[(-1)^n+1]n$, odd terms (orange) equal $0$, while even terms (blue) equal $2n$ and tend to $+\infty$.}
  \label{fig:limsup-extended-subsequences}
\end{illustrationfigure}
Figure \ref{fig:limsup-extended-subsequences} exhibits finite and infinite subsequential limits in one sequence.

\begin{definition} \label{def:upper-and-lower-limits} \hlt{Upper and Lower Limits}\\
Let $(s_n)$ be a real sequence and set $E\defeq E(s)\defeq\{x\in\overline{\R}:(\exists\,n_1<n_2<\cdots)\ s_{n_k}\to x\}$.
\begin{enumerate}[label=(\roman*)]
\setlength{\itemsep}{0pt}
\item The \hlt{upper limit} is $\limsup\limits_{n\to\infty}s_n\defeq s^*\defeq\sup\limits_{\overline{\R}}E$.
\item The \hlt{lower limit} is $\liminf\limits_{n\to\infty}s_n\defeq s_*\defeq\inf\limits_{\overline{\R}}E$.
\end{enumerate}
\end{definition}

\begin{theorem} \label{thm:limsup-characterization} \hlt{Characterization of Upper and Lower Limits}\\
Let $(s_n)$ be a real sequence and set $E\defeq\{x\in\overline{\R}:(\exists\,n_1<n_2<\cdots)\ s_{n_k}\to x\}$. Set $s^*\defeq\sup\limits_{\overline{\R}}E$ and $s_*\defeq\inf\limits_{\overline{\R}}E$.\\
For $y\in\overline{\R}$, set $U(y)\defeq(\forall x\in\R)[x>y\Rightarrow(\exists N\in\N)(\forall n\geq N)\,s_n<x]$. Set $L(y)\defeq(\forall x\in\R)[x<y\Rightarrow(\exists N\in\N)(\forall n\geq N)\,s_n>x]$.
\begin{enumerate}[label=(\roman*)]
\setlength{\itemsep}{0pt}
\item $s^*\in E$.
\item $U(s^*)$.
\item $\{y\in E:U(y)\}=\{s^*\}$.
\item $s_*\in E$.
\item $L(s_*)$.
\item $\{y\in E:L(y)\}=\{s_*\}$.
\end{enumerate}
\begin{proof}
\begin{enumerate}[label=(\roman*)]
\setlength{\itemsep}{0pt}
\item Set $n_0\defeq0$ and distinguish the value of $s^*$.
\begin{enumerate}[label=(\alph*)]
\setlength{\itemsep}{0pt}
\item If $s^*=+\infty$, choose $y_k\in E$ with $y_k>k$. Choose $n_k>n_{k-1}$ with $s_{n_k}>k$; this is possible by $y_k\in E$ and $y_k>k$.\\
For $M\in\R$, choose $K\in\N$ with $K>M$ by Theorem \ref{thm:archimedean-and-rational-density}(i). Then $(\forall k\geq K)\,s_{n_k}>k>M$, so $s_{n_k}\to+\infty$ and $+\infty\in E$.
\item If $s^*\in\R$, choose $y_k\in E\cap(s^*-k^{-1},s^*]$ and then $n_k>n_{k-1}$ with $\abs{s_{n_k}-y_k}<k^{-1}$.\\
The bound $\abs{s_{n_k}-s^*}<2k^{-1}$ holds by Proposition \ref{prop:real-absolute-value-triangle-inequality}(iv). For $\varepsilon>0$, choose $K>2\varepsilon^{-1}$ by Theorem \ref{thm:archimedean-and-rational-density}(i).\\
Then $(\forall k\geq K)\,\abs{s_{n_k}-s^*}<\varepsilon$, so $s_{n_k}\to s^*$ and $s^*\in E$.
\item Suppose $s^*=-\infty$ and $-\infty\notin E$. Some tail is bounded below.\\
Otherwise the recursive construction with $s_{n_k}\leq-k$ would place $-\infty$ in $E$.\\
If this tail is bounded above, it has a finite subsequential limit by Theorem \ref{thm:bolzano-weierstrass-real-sequence}.\\
If it is unbounded above, it has an infinite subsequential limit by Proposition \ref{prop:unbounded-real-sequence-has-infinite-subsequence}.\\
The infinite limit is $+\infty$ because the tail is bounded below.\\
Either case places an element greater than $-\infty$ in $E$, contradicting $s^*=\sup\limits_{\overline{\R}}E=-\infty$.
\end{enumerate}
\item Fix $x\in\R$ with $x>s^*$. If $s_n\geq x$ infinitely often, choose a subsequence with $(\forall k\in\N)\,s_{n_k}\geq x$.\\
If subsequence bounded, a further subsequence converges by Theorem \ref{thm:bolzano-weierstrass-real-sequence}. Its limit $y$ satisfies $y\geq x$ by Proposition \ref{prop:convergent-sequence-basic-consequences}(iii).\\
If the subsequence is unbounded, it is bounded below by $x$. A further subsequence tends to $+\infty$ by Proposition \ref{prop:unbounded-real-sequence-has-infinite-subsequence}.\\
Thus $E\cap[x,+\infty]\neq\emptyset$, contradicting $x>s^*=\sup\limits_{\overline{\R}}E$. Therefore $s_n<x$ eventually.
\item Let $A\defeq\{y\in E:U(y)\}$. Then $s^*\in A$ by (i)--(ii). If $p,q\in A$ with $p<q$, choose $x\in\R$ with $p<x<q$.\\
Then $s_n<x$ eventually by $U(p)$. Since $q\in E$, $q>x$, a subsequence is eventually greater than $x$, contradiction. Thus $A=\{s^*\}$.
\item For $r_n\defeq-s_n$, its subsequential-limit set is $-E$, upper limit is $-s_*$.\\
Relation $-s_*\in-E$ holds by (i) for $(r_n)$; equivalently, $s_*\in E$.
\item Fix $x\in\R$ with $x<s_*$. Then $-x>-s_*$, so $-s_n<-x$ eventually by (ii) for $(r_n)$. Equivalently, $s_n>x$ eventually.
\item The set $\{y\in E:L(y)\}$ is the negative of the upper-characterizing set for $(r_n)$, so it equals $\{s_*\}$ by (iii).
\end{enumerate}
\hlt{Structural Commentary}: Diagonal subsequences attain each extremum; sign reversal transfers the upper characterization downward.\\
\hlt{Key Insight}: Failure of an eventual bound creates a competing subsequential limit.
\end{proof}
\end{theorem}

\begin{proposition} \label{prop:limsup-liminf-tail-suprema} \hlt{Tail-Envelope Formulae}\\
For a real sequence $(s_n)$ and $N\in\N$, set $T_N\defeq\{s_n:n\geq N\}$. Set $u_N\defeq\sup\limits_{\overline{\R}}T_N$ and $\ell_N\defeq\inf\limits_{\overline{\R}}T_N$.
\begin{enumerate}[label=(\roman*)]
\setlength{\itemsep}{0pt}
\item $(\forall N\in\N)\,[u_{N+1}\leq u_N\land\ell_{N+1}\geq\ell_N]$.
\item $\limsup\limits_{n\to\infty}s_n=\inf\limits_{N\in\N}u_N$ and $\liminf\limits_{n\to\infty}s_n=\sup\limits_{N\in\N}\ell_N$.
\end{enumerate}
\begin{proof}
\begin{enumerate}[label=(\roman*)]
\setlength{\itemsep}{0pt}
\item The inclusion $T_{N+1}\subseteq T_N$ implies $u_{N+1}\leq u_N$ and $\ell_{N+1}\geq\ell_N$.
\item Set $s^*\defeq\limsup s_n$ and $u\defeq\inf\limits_{N\in\N}u_N$. If $x\in\R$ and $x>s^*$, then $u\leq x$ by Theorem \ref{thm:limsup-characterization}(ii) and the definition of $u$.\\
If $y\in\R$ and $y<s^*$, each tail contains a term greater than $y$ by Theorem \ref{thm:limsup-characterization}(i). Thus $(\forall N\in\N)\,u_N\geq y$ and $u\geq y$.\\
If $s^*=-\infty$, then $(\forall x\in\R)\,u\leq x$, so $u=-\infty$. If $s^*=+\infty$, then $(\forall y\in\R)\,u\geq y$, so $u=+\infty$.\\
If $s^*\in\R$, then $(\forall\varepsilon>0)\,s^*-\varepsilon\leq u\leq s^*+\varepsilon$, so $u=s^*$. For $r_n\defeq-s_n$, its upper tail extrema are $-\ell_N$ and $\limsup r_n=-s_*$.\\
Thus $-s_*=\inf\limits_{N\in\N}(-\ell_N)=-\sup\limits_{N\in\N}\ell_N$, so $\sup\limits_{N\in\N}\ell_N=s_*$.
\end{enumerate}
\hlt{Structural Commentary}: Nested tails make extrema monotone; subsequential thresholds identify their terminal bounds.\\
\hlt{Key Insight}: Every extremal subsequence returns arbitrarily far into each tail.
\end{proof}
\end{proposition}

\begin{illustrationfigure}[H]
  \centering
  \begin{tikzpicture}[
  >=Stealth,
  x=0.9cm,
  y=0.9cm,
  every node/.style={font=\small},
  axis/.style={->, gray!70, line width=0.55pt},
  upper/.style={blue!70!black, line width=0.75pt},
  lower/.style={orange!85!black, line width=0.75pt},
  udot/.style={circle, fill=blue!70!black, inner sep=1.25pt},
  ldot/.style={circle, fill=orange!85!black, inner sep=1.25pt},
  limit/.style={green!45!black, dashed, line width=0.55pt}
]
  \draw[axis] (0.55,0) -- (7.55,0) node[right] {$N$};
  \draw[axis] (0.85,-1.55) -- (0.85,1.65);

  \draw[limit] (1.00,0.62) -- (7.20,0.62);
  \draw[limit] (1.00,-0.62) -- (7.20,-0.62);
  \node[green!45!black, anchor=west] at (7.25,0.62) {$s^*$};
  \node[green!45!black, anchor=west] at (7.25,-0.62) {$s_*$};

  \coordinate (u1) at (1.25,1.34);
  \coordinate (u2) at (2.15,1.04);
  \coordinate (u3) at (3.05,0.86);
  \coordinate (u4) at (3.95,0.75);
  \coordinate (u5) at (4.85,0.68);
  \coordinate (u6) at (5.75,0.64);
  \coordinate (u7) at (6.65,0.625);
  \draw[upper] plot[smooth] coordinates {(u1) (u2) (u3) (u4) (u5) (u6) (u7)};
  \node[blue!70!black, anchor=south] at (2.15,1.12) {$u_N$};

  \coordinate (l1) at (1.25,-1.34);
  \coordinate (l2) at (2.15,-1.03);
  \coordinate (l3) at (3.05,-0.85);
  \coordinate (l4) at (3.95,-0.74);
  \coordinate (l5) at (4.85,-0.68);
  \coordinate (l6) at (5.75,-0.64);
  \coordinate (l7) at (6.65,-0.625);
  \draw[lower] plot[smooth] coordinates {(l1) (l2) (l3) (l4) (l5) (l6) (l7)};
  \node[orange!85!black, anchor=north] at (2.15,-1.11) {$\ell_N$};

  \foreach \p in {u1,u2,u3,u4,u5,u6,u7} {
    \node[udot] at (\p) {};
  }
  \foreach \p in {l1,l2,l3,l4,l5,l6,l7} {
    \node[ldot] at (\p) {};
  }
\end{tikzpicture}

  \caption{The upper tail extrema $u_N$ decrease to $s^*=\limsup s_n$, while the lower extrema $\ell_N$ increase to $s_*=\liminf s_n$.}
  \label{fig:limsup-liminf-tail-bounds}
\end{illustrationfigure}
Figure \ref{fig:limsup-liminf-tail-bounds} visualizes the two monotone envelopes in Proposition \ref{prop:limsup-liminf-tail-suprema}.

\begin{theorem} \label{thm:convergence-iff-limsup-liminf-equal} \hlt{Convergence and Equal Upper/Lower Limits}\\
For a real sequence $(s_n)$ and $s\in\R$, $s_n\to s\Longleftrightarrow\limsup\limits_{n\to\infty}s_n=\liminf\limits_{n\to\infty}s_n=s$.
\begin{proof}
\begin{enumerate}[label=\roman*.]
\setlength{\itemsep}{0pt}
\item[$(\Rightarrow)$] Suppose $s_n\to s$ and set $E\defeq E(s)$. Every subsequence converges to $s$ by Proposition \ref{prop:subsequences-of-convergent-sequences}(i).\\
Choose $M>0$, $(\forall n\in\N)\,\abs{s_n}\leq M$ by Proposition \ref{prop:convergent-sequence-basic-consequences}(ii). Then $(\forall n\in\N)\,-M\leq s_n\leq M$, no subsequence tends to $\pm\infty$.\\
Thus $E=\{s\}$ and $\limsup s_n=\liminf s_n=s$.
\item[$(\Leftarrow)$] Suppose $\limsup s_n=\liminf s_n=s$, $\varepsilon>0$. Inequalities $s_n<s+\varepsilon$ and $s_n>s-\varepsilon$ hold eventually by Theorem \ref{thm:limsup-characterization}(ii),(v).\\
Therefore $\abs{s_n-s}<\varepsilon$ eventually, so $s_n\to s$.
\end{enumerate}
\hlt{Structural Commentary}: Subsequence inheritance proves necessity; the two eventual extremal bounds prove sufficiency.\\
\hlt{Key Insight}: Limsup supplies an eventual ceiling; liminf supplies an eventual floor.
\end{proof}
\end{theorem}

\begin{theorem} \label{thm:limsup-liminf-order-comparison} \hlt{Order Comparison for Upper and Lower Limits}\\
Let $(s_n)$ and $(t_n)$ be real sequences with $s_n\leq t_n$ eventually.
\begin{enumerate}[label=(\roman*)]
\setlength{\itemsep}{0pt}
\item $\limsup\limits_{n\to\infty}s_n\leq \limsup\limits_{n\to\infty}t_n$.
\item $\liminf\limits_{n\to\infty}s_n\leq \liminf\limits_{n\to\infty}t_n$.
\end{enumerate}
\begin{proof}
Choose $N\in\N$ such that $(\forall n\geq N)\,s_n\leq t_n$ by the definition of eventual order.
\begin{enumerate}[label=(\roman*)]
\setlength{\itemsep}{0pt}
\item $(\forall k\geq N)\,\sup\limits_{n\geq k}s_n\leq\sup\limits_{n\geq k}t_n$. Then $\limsup s_n=\inf\limits_{k\geq N}\sup\limits_{n\geq k}s_n\leq\inf\limits_{k\geq N}\sup\limits_{n\geq k}t_n=\limsup t_n$ by Proposition \ref{prop:limsup-liminf-tail-suprema}(i)--(ii).
\item $(\forall k\geq N)\,\inf\limits_{n\geq k}s_n\leq\inf\limits_{n\geq k}t_n$. Then $\liminf s_n=\sup\limits_{k\geq N}\inf\limits_{n\geq k}s_n\leq\sup\limits_{k\geq N}\inf\limits_{n\geq k}t_n=\liminf t_n$ by Proposition \ref{prop:limsup-liminf-tail-suprema}(i)--(ii).
\end{enumerate}
\hlt{Structural Commentary}: Tailwise extrema transfer eventual pointwise order to both envelope limits.\\
\hlt{Key Insight}: Deleting finitely many tail extrema preserves their terminal infimum and supremum.
\end{proof}
\end{theorem}

\begin{theorem} \label{thm:pinching-theorem-sequences} \hlt{Pinching Theorem for Real Sequences}\\
Let $(a_n)$, $(s_n)$, and $(b_n)$ be real sequences, and let $s\in\R$. If $a_n\leq s_n\leq b_n$ eventually and $a_n,b_n\to s$, then $s_n\to s$.
\begin{proof}
Fix $\varepsilon>0$ and choose $N_0$ such that $(\forall n\geq N_0)\,a_n\leq s_n\leq b_n$. Choose $N_1,N_2$ such that $(\forall n\geq N_1)\,s-\varepsilon<a_n$ and $(\forall n\geq N_2)\,b_n<s+\varepsilon$.\\
Set $N\defeq\max\{N_0,N_1,N_2\}$. Then $(\forall n\geq N)\,s-\varepsilon<s_n<s+\varepsilon$, so $s_n\to s$.\\
\hlt{Structural Commentary}: Three eventual thresholds combine at their maximum to trap the middle sequence.\\
\hlt{Key Insight}: $N=\max\{N_0,N_1,N_2\}$ synchronizes both convergent bounds.
\end{proof}
\end{theorem}

\begin{illustrationfigure}[H]
  \centering
  \begin{tikzpicture}[
  >=Stealth,
  x=0.9cm,
  y=0.9cm,
  every node/.style={font=\small},
  axis/.style={->, gray!70, line width=0.55pt},
  upper/.style={blue!70!black, line width=0.9pt},
  lower/.style={orange!85!black, line width=0.9pt},
  upperdot/.style={circle, fill=blue!70!black, inner sep=1.35pt},
  lowerdot/.style={circle, fill=orange!85!black, inner sep=1.35pt},
  middledot/.style={circle, draw=green!45!black, fill=white, line width=0.65pt, inner sep=1.45pt}
]
  \draw[axis] (0.35,0) -- (7.55,0) node[right] {$n$};
  \draw[axis] (0.7,-1.45) -- (0.7,1.55);
  \node[left] at (0.65,0) {$s$};

  \coordinate (u1) at (1,1.18);
  \coordinate (u2) at (2,0.72);
  \coordinate (u3) at (3,0.45);
  \coordinate (u4) at (4,0.28);
  \coordinate (u5) at (5,0.18);
  \coordinate (u6) at (6,0.10);
  \coordinate (u7) at (7,0.05);

  \draw[upper]
    plot[smooth] coordinates {(u1) (u2) (u3) (u4) (u5) (u6) (u7)};

  \coordinate (l1) at (1,-1.08);
  \coordinate (l2) at (2,-0.60);
  \coordinate (l3) at (3,-0.34);
  \coordinate (l4) at (4,-0.20);
  \coordinate (l5) at (5,-0.11);
  \coordinate (l6) at (6,-0.06);
  \coordinate (l7) at (7,-0.03);

  \draw[lower]
    plot[smooth] coordinates {(l1) (l2) (l3) (l4) (l5) (l6) (l7)};

  \foreach \p in {u1,u2,u3,u4,u5,u6,u7} {
    \node[upperdot] at (\p) {};
  }
  \foreach \p in {l1,l2,l3,l4,l5,l6,l7} {
    \node[lowerdot] at (\p) {};
  }

  \foreach \x/\ym in {
    1/-0.42,
    2/0.16,
    3/0.02,
    4/0.12,
    5/-0.02,
    6/0.03,
    7/0.01
  } {
    \node[middledot] at (\x,\ym) {};
  }

  \node[blue!70!black, anchor=west] at (1.08,1.20) {$(b_n)$};
  \node[green!45!black, anchor=west] at (1.08,-0.42) {$(s_n)$};
  \node[orange!85!black, anchor=west] at (1.08,-1.2) {$(a_n)$};
\end{tikzpicture}

  \caption{With $a_n\leq s_n\leq b_n$ and $a_n,b_n\to s$, the middle sequence also tends to $s$.}
  \label{fig:pinching-theorem-sequences}
\end{illustrationfigure}
Figure \ref{fig:pinching-theorem-sequences} shows the middle sequence trapped between the two convergent bounds.

\begin{theorem} \label{thm:special-sequence-limits} \hlt{Special Sequence Limits}
\begin{enumerate}[label=(\roman*)]
\setlength{\itemsep}{0pt}
\item If $p>0$, then $\lim\limits_{n\to\infty}n^{-p}=0$.
\item If $p>0$, then $\lim\limits_{n\to\infty}\sqrt[n]{p}=1$.
\item $\lim\limits_{n\to\infty}\sqrt[n]{n}=1$.
\item If $p>0$ and $\alpha\in\R$, then $\lim\limits_{n\to\infty}\frac{n^\alpha}{(1+p)^n}=0$.
\item If $\abs{x}<1$, then $\lim\limits_{n\to\infty}x^n=0$.
\end{enumerate}
\begin{proof} All positive roots below exist uniquely by Theorem \ref{thm:nth-roots-positive-reals}.
\begin{enumerate}[label=(\roman*)]
\setlength{\itemsep}{0pt}
\item Fix $\varepsilon>0$ and choose $N>\varepsilon^{-p^{-1}}$ by Theorem \ref{thm:archimedean-and-rational-density}(i). Then $(\forall n\geq N)\,0<n^{-p}<\varepsilon$.
\item Suppose $p>1$ and set $x_n\defeq\sqrt[n]{p}-1$. Then $x_n>0$ and $1+nx_n\leq(1+x_n)^n=p$ by Proposition \ref{prop:binomial-theorem-and-subset-count}(i), so $0<x_n\leq(p-1)n^{-1}$.\\
Upper bound tends to $0$ by (i) and Theorem \ref{thm:complex-sequence-limit-algebra}(ii). Thus $x_n\to0$ by Theorem \ref{thm:pinching-theorem-sequences}; $\sqrt[n]{p}=1+x_n\to1$ by Theorem \ref{thm:complex-sequence-limit-algebra}(ii).\\
If $p=1$, then $(\forall n\in\N)\,\sqrt[n]{p}=1$. If $0<p<1$, set $q\defeq p^{-1}>1$. Then $\sqrt[n]{q}\to1$ by the preceding case.\\
The identity $\sqrt[n]{p}=(\sqrt[n]{q})^{-1}$ holds by Corollary \ref{cor:nth-root-product-positive-reals}. Therefore $\sqrt[n]{p}\to1$ by Theorem \ref{thm:complex-sequence-limit-algebra}(iv).
\item Set $x_n\defeq\sqrt[n]{n}-1\geq0$. For $n\geq2$, $n=(1+x_n)^n\geq\binom{n}{2}x_n^2=\frac{n(n-1)}{2}x_n^2$ by Proposition \ref{prop:binomial-theorem-and-subset-count}(i). Thus $0\leq x_n^2\leq\frac{2}{n-1}$.\\
Fix $\varepsilon>0$. Choose $N>1+2\varepsilon^{-2}$ by Theorem \ref{thm:archimedean-and-rational-density}(i). For $n\geq\max\{2,N\}$, $0\leq x_n<\varepsilon$. Hence $x_n\to0$ and $\sqrt[n]{n}=1+x_n\to1$.
\item Choose $k\in\N$ with $k>\alpha$ by Theorem \ref{thm:archimedean-and-rational-density}(i). If $n>2k$, then $(\forall j\in\{0,\ldots,k-1\})\,n-j>\frac{n}{2}$.\\
Thus  $(1+p)^n\geq\binom{n}{k}p^k>\frac{n^kp^k}{2^kk!}$ by Proposition \ref{prop:binomial-theorem-and-subset-count}(i); $0<\frac{n^\alpha}{(1+p)^n}<\frac{2^kk!}{p^k}n^{-(k-\alpha)}$.\\
Upper bound tends to $0$ by (i), Theorem \ref{thm:complex-sequence-limit-algebra}(ii). Quotient tends to $0$ by Theorem \ref{thm:pinching-theorem-sequences}.
\item If $x\neq0$, set $p\defeq\abs{x}^{-1}-1>0$. Then $\abs{x}=(1+p)^{-1}$ and $\abs{x}^n\to0$ by (iv) with $\alpha=0$.\\
The bounds $-\abs{x}^n\leq x^n\leq\abs{x}^n$ imply $x^n\to0$ by Theorem \ref{thm:pinching-theorem-sequences}. The case $x=0$ is immediate.
\end{enumerate}
\hlt{Structural Commentary}: Reciprocal powers and binomial lower bounds reduce roots and exponential quotients to pinching.\\
\hlt{Key Insight}: A fixed binomial term dominates every lower polynomial power.
\end{proof}
\end{theorem}

\begin{proposition} \label{prop:decimal-truncations-squeeze-real} \hlt{Decimal Truncations Squeeze a Real}\\
Let $x\geq0$ and set $S_n\defeq\{k\in\Z_{\geq0}:k10^{-n}\leq x\}$ for $n\in\N_0$.
\begin{enumerate}[label=(\roman*)]
\setlength{\itemsep}{0pt}
\item The maximum $s_n\defeq\max S_n$ exists and satisfies $s_n10^{-n}\leq x<(s_n+1)10^{-n}$ for $n\in\N_0$.
\item $(\forall n\in\N_0)\,0\leq x-s_n10^{-n}<10^{-n}$, and $s_n10^{-n}\to x$.
\item $(\forall n\in\N_0)(\exists!a_{n+1}\in\{0,\ldots,9\})\,s_{n+1}=10s_n+a_{n+1}$.
\end{enumerate}
\begin{proof}
\begin{enumerate}[label=(\roman*)]
\setlength{\itemsep}{0pt}
\item Fix $n\in\N_0$ and set $A_n\defeq\{j\in\N:j10^{-n}>x\}$. The set $A_n$ is nonempty by Theorem \ref{thm:archimedean-and-rational-density}(i).\\
Let $m_n$ be its least element by the standing well-ordering assumption.\\
Since $x\geq0$, $m_n\geq1$, and minimality implies $m_n-1\in S_n$. Every $k\in S_n$ satisfies $k<m_n$, so $s_n=m_n-1=\max S_n$.\\
The inequalities $s_n10^{-n}\leq x<(s_n+1)10^{-n}$ follow from $s_n\in S_n$ and $s_n+1=m_n\in A_n$.
\item The inequalities in (i) imply $(\forall n\in\N_0)\,0\leq x-s_n10^{-n}<10^{-n}$. The limit $10^{-n}\to0$ holds by Theorem \ref{thm:special-sequence-limits}(v) with $x=10^{-1}$.\\
Thus $x-10^{-n}\to x$ by Theorem \ref{thm:complex-sequence-limit-algebra}(ii). The bounds $x-10^{-n}\leq s_n10^{-n}\leq x$ imply $s_n10^{-n}\to x$ by Theorem \ref{thm:pinching-theorem-sequences}.
\item The relations $10s_n\in S_{n+1}$ and $10s_n+10\in A_{n+1}$ imply $10s_n\leq s_{n+1}<10s_n+10$.\\
Thus $a_{n+1}\defeq s_{n+1}-10s_n\in\{0,\ldots,9\}$, and this formula makes $a_{n+1}$ unique.
\end{enumerate}
\hlt{Structural Commentary}: Archimedean grids define each truncation; nested mesh bounds then support the squeeze.\\
\hlt{Key Insight}: Refinement by a factor of $10$ leaves one digit as the integer remainder.
\end{proof}
\end{proposition}

\begin{illustrationfigure}[H]
  \centering
  \begin{tikzpicture}[
  >=Stealth,
  x=0.9cm,
  y=0.9cm,
  every node/.style={font=\small},
  axis/.style={gray!70, line width=0.55pt},
  cell/.style={blue!70!black, line width=1.0pt},
  closedendpoint/.style={circle, fill=blue!70!black, inner sep=1.45pt},
  openendpoint/.style={circle, draw=blue!70!black, fill=white, line width=0.75pt, inner sep=1.45pt},
  xpoint/.style={circle, draw=black!75, fill=white, line width=0.6pt, inner sep=1.55pt},
  widthmark/.style={orange!75!black, line width=0.65pt}
]
  \draw[axis] (-0.45,0) -- (12.45,0);

  \coordinate (L) at (1.55,0);
  \coordinate (X) at (6.0,0);
  \coordinate (R) at (10.45,0);
  \coordinate (LW) at (1.55,0.48);
  \coordinate (RW) at (10.45,0.48);

  \draw[cell] (L) -- (R);
  \node[closedendpoint] at (L) {};
  \node[xpoint] at (X) {};
  \node[openendpoint] at (R) {};

  \node[anchor=base] at (1.55,-0.45) {$s_n10^{-n}$};
  \node at (6.0,0.27) {$x$};
  \node[anchor=base] at (10.45,-0.45) {$(s_n+1)10^{-n}$};

  \node[anchor=base] at (6.0,-0.45) {$s_n10^{-n}\le x<(s_n+1)10^{-n}$};

  \draw[widthmark] (LW) -- (RW);
  \draw[widthmark] (LW) -- ++(0,0.18);
  \draw[widthmark] (RW) -- ++(0,0.18);
  \node[anchor=base, font=\scriptsize, orange!75!black] at (6.0,0.75) {cell width $10^{-n}$};
\end{tikzpicture}

  \caption{At precision $10^{-n}$, $x$ lies in the half-open cell $[s_n10^{-n},(s_n+1)10^{-n})$ of width $10^{-n}$.}
  \label{fig:decimal-truncation-squeeze}
\end{illustrationfigure}
Figure \ref{fig:decimal-truncation-squeeze} illustrates the shrinking interval in Proposition \ref{prop:decimal-truncations-squeeze-real}(i).

\begin{proofmethod} \label{meth:sequence-proof-patterns} \hlt{Real-Sequence Proof Patterns}
\begin{enumerate}[label=(\roman*)]
\setlength{\itemsep}{0pt}
\item \hlt{Convergence}: Fix $\varepsilon>0$ and construct $N$ such that $(\forall n\geq N)\,\abs{s_n-s}<\varepsilon$.
\item \hlt{Uniqueness}: Synchronize two $\frac{\varepsilon}{2}$ tail bounds and compare the proposed limits by the triangle inequality.
\item \hlt{Subsequences}: Recursively choose $n_k>n_{k-1}$ so the $k$th term meets a threshold or an error bound of $k^{-1}$.
\item \hlt{Cauchy transfer}: Anchor an arbitrary late term to a sufficiently late convergent subsequence term.
\item \hlt{Monotone convergence}: Choose the relevant supremum or infimum and propagate one extremal witness through its tail.
\end{enumerate}
\hlt{Key Insight}: Synchronized tail indices reduce global limit claims to one controlling estimate.
\end{proofmethod}
